% 00-frontmatter.tex — Trang bìa, lời nói đầu, tóm tắt điều hành, danh mục từ viết tắt

% ===================== TRANG BÌA =====================
\begin{titlepage}
\centering
\vspace*{0.8cm}

{\large\bfseries TỔNG CÔNG TY BƯU ĐIỆN VIỆT NAM}\\[0.25cm]
{\normalsize TRUNG TÂM ĐỔI MỚI SÁNG TẠO – PHÒNG AI}\\[1.5cm]

{\color{navy}\rule{\textwidth}{0.8pt}}\\[0.9cm]

\begin{minipage}{0.86\textwidth}
\centering
{\bfseries\large
Triển khai các mô hình dự báo, dự đoán về hành vi của khách hàng
phục vụ công tác cải thiện chất lượng dịch vụ và giữ chân người dùng.\par}
\end{minipage}\\[0.9cm]

{\color{navy}\rule{\textwidth}{0.8pt}}\\[0.8cm]

{\bfseries BÁO CÁO KHOA HỌC}\\[2.6cm]

\begin{tabular}{@{}r@{\hspace{8mm}}l@{}}
\textbf{Kính gửi:} & Ban lãnh đạo Tổng công ty Bưu điện Việt Nam \\[0.22cm]
\textbf{Đơn vị thực hiện:} & Phòng AI, Trung tâm Đổi mới Sáng tạo \\[0.22cm]
\textbf{Ngày báo cáo:} & 13/10/2026 \\
\end{tabular}

\vfill
{\large\bfseries 2026}
\end{titlepage}

% ===================== LỜI NÓI ĐẦU =====================
\chapter*{Lời nói đầu}
\addcontentsline{toc}{chapter}{Lời nói đầu}
Trong hoạt động bưu chính - chuyển phát, khách hàng không nói với doanh nghiệp mọi điều
họ nghĩ, nhưng họ để lại dấu vết trong dữ liệu vận hành: gửi bao nhiêu đơn, vào những ngày
nào, đơn giao có đúng hẹn không, có khiếu nại không, doanh thu mỗi tháng tăng hay giảm.
Khoa học dữ liệu cho phép khai thác các dấu vết đó một cách có hệ thống để hiểu khách hàng
đang làm gì, vì sao, và nhiều khả năng sẽ làm gì tiếp theo. Báo cáo khoa học này trình bày
kết quả nghiên cứu và triển khai của đề tài ``Triển khai các mô hình dự báo, dự đoán về
hành vi của khách hàng phục vụ công tác cải thiện chất lượng dịch vụ và giữ chân người
dùng'' tại Tổng công ty Bưu điện Việt Nam, tập trung vào bài toán dự báo khách hàng có nguy cơ rời bỏ dịch vụ.

Báo cáo được tổ chức theo trình tự từ vấn đề đến giải pháp và kiểm nghiệm: giới thiệu đề
tài và bối cảnh; mục tiêu và phạm vi; nội dung giải pháp; cách thức
triển khai và vận hành tự động; nguồn lực và tiến độ; kết quả áp dụng; hiệu quả và bộ chỉ
tiêu đo lường; đánh giá và quản lý rủi ro; cùng khả năng áp dụng, mở rộng cho các bài
toán hành vi khác.

Toàn bộ nội dung trình bày trong báo cáo được xây dựng từ dữ liệu vận hành có sẵn của Tổng
công ty. Đây là báo cáo tiến độ: hệ thống chưa vận hành chính thức nên các số liệu đánh giá
mô hình và hiệu quả kinh doanh chưa được đưa vào, sẽ bổ sung ở báo cáo sau.

Xin gửi lời cảm ơn tới Ban lãnh đạo và các đơn vị nghiệp vụ, kỹ thuật dữ liệu đã tạo điều
kiện, chia sẻ bối cảnh nghiệp vụ và dữ liệu để hoàn thành báo cáo này.

\vspace{1cm}
\begin{flushright}
\emph{Phòng AI, Trung tâm Đổi mới Sáng tạo}
\end{flushright}

% ===================== TÓM TẮT ĐIỀU HÀNH =====================
\chapter*{Tóm tắt điều hành}
\addcontentsline{toc}{chapter}{Tóm tắt điều hành}
Tổng công ty Bưu điện Việt Nam phục vụ một lượng lớn khách hàng bưu chính, chuyển phát. Quan
trọng nhất là nhóm khách hàng doanh nghiệp, cửa hàng và đơn vị thương mại điện tử gửi hàng
thường xuyên theo hợp đồng: họ mang lại sản lượng và doanh thu ổn định, nhưng cũng có nhiều
nhà cung cấp khác để lựa chọn. Với nhóm này, giữ được một khách đang dùng dịch vụ thường có
giá trị hơn việc tìm một khách mới tương đương, và muốn giữ được thì phải phát hiện sớm dấu
hiệu khách bắt đầu thay đổi thói quen. Báo cáo này nghiên cứu cách dùng dữ liệu vận hành sẵn
có để dự báo hành vi khách hàng, tập trung vào bài toán cảnh báo sớm khách có nguy
cơ rời bỏ. Rời bỏ được hiểu rộng, gồm cả ngừng hẳn, giảm mạnh số đơn, giảm mạnh doanh thu
hoặc đi xuống rõ rệt.

Giải pháp là một dây chuyền sáu bước đi từ dữ liệu thô đến danh sách khách cần chăm sóc. Hệ
thống chia khách thành các nhóm và giữ lại những nhóm có đủ dữ liệu để dự báo, rồi tính các
chỉ số hành vi bằng cách so sánh mỗi khách với mức quen thuộc của chính họ. Để máy học được
thế nào là rời bỏ, các khách được gán nhãn bằng bốn quy tắc nghiệp vụ dựa trên đơn hàng
thực tế, được xây dựng có tham khảo ý kiến của đội chăm sóc khách hàng. Mô hình gồm hai lớp, một mô hình cơ sở làm mốc so sánh và một mô hình chính
thuộc họ cây quyết định tăng cường, được kiểm tra theo trình tự thời gian như khi
dùng thật: học từ các tháng cũ, kiểm tra trên tháng mới hơn và giữ tháng gần nhất làm bài kiểm tra. Chất lượng được đo bằng các thước đo gắn với nghiệp vụ như F1, độ chính xác của
cảnh báo và độ bao phủ, cùng PR-AUC và ROC-AUC. Mỗi dự báo đi kèm tối đa ba lý do rủi ro
thuộc tám nhóm nghiệp vụ, rút ra bằng kỹ thuật phân tích SHAP, để nhân viên chăm sóc biết
cần hỏi khách điều gì.

Toàn bộ hệ thống chạy tự động qua sáu luồng công việc nối tiếp trên apache airflow: nạp và
kiểm tra dữ liệu hằng ngày, tạo chỉ số và gán nhãn, huấn luyện lại theo chu kỳ và chỉ thay mô
hình khi mô hình mới tốt hơn và qua kiểm định, chấm điểm và xuất danh sách, giám sát độ lệch
của dữ liệu và điểm số bằng PSI và KS, cùng dọn dẹp định kỳ. Hệ thống dùng hạ tầng sẵn có của
đơn vị, lưu dữ liệu bằng PostgreSQL, đóng gói bằng container và có quy trình phát hành, hoàn
tác phiên bản. Mỗi kỳ, danh sách ưu tiên chỉ gồm 5\% khách có rủi ro cao nhất, đủ ít để đội
chăm sóc xử lý được thực sự.

Dây chuyền từ dữ liệu thô đến danh sách kèm lý do đã được xây dựng hoàn chỉnh về mặt kỹ thuật
nhưng chưa vận hành chính thức, nên các số liệu về chất lượng mô hình và hiệu quả kinh doanh
sẽ được bổ sung ở báo cáo sau.
Báo cáo cũng nêu rõ những hạng mục chưa hoàn thành cùng cách kiểm soát rủi ro tương ứng, và đề xuất lộ
trình nâng cấp cũng như cách tái sử dụng nền tảng này cho các bài toán hành vi khác.

\vspace{0.4cm}
\noindent\textbf{Từ khóa:} dự báo hành vi khách hàng, rời bỏ dịch vụ, giữ chân khách
hàng, học máy, chỉ số hành vi, kiểm tra theo thời gian, SHAP, apache airflow.

% ===================== DANH MỤC TỪ VIẾT TẮT =====================
\chapter*{Danh mục từ viết tắt}
\addcontentsline{toc}{chapter}{Danh mục từ viết tắt}
\begin{longtable}{@{}p{3.2cm}p{11.5cm}@{}}
\toprule
\textbf{Từ viết tắt} & \textbf{Diễn giải} \\
\midrule
\endfirsthead
\toprule
\textbf{Từ viết tắt} & \textbf{Diễn giải} \\
\midrule
\endhead
\bottomrule
\endfoot
Airflow & Apache Airflow - phần mềm điều phối luồng công việc theo lịch \\
API & Application Programming Interface - giao diện lập trình ứng dụng \\
AUC & Area Under the Curve - diện tích dưới đường cong của thước đo xếp hạng \\
COD & Cash on Delivery -- thu hộ tiền khi giao hàng \\
CSV & Comma-Separated Values - tệp dữ liệu phân tách bằng dấu phẩy \\
DAG & Directed Acyclic Graph - đồ thị có hướng không chu trình, khai báo luồng công việc \\
F1 & Điểm F1 - cân bằng giữa độ chính xác cảnh báo và độ bao phủ \\
KS & Kolmogorov Smirnov - phép đo độ lệch phân phối dữ liệu \\
My VNPost & Ứng dụng bưu gửi của Tổng công ty Bưu điện Việt Nam \\
PR & Precision Recall - đường cong độ chính xác/bao phủ \\
PSI & Population Stability Index - chỉ số ổn định dân số, đo độ lệch dữ liệu \\
ROC & Receiver Operating Characteristic - đường cong đặc trưng hoạt động \\
SHAP & SHapley Additive exPlanations - mức đóng góp của từng chỉ số vào kết quả dự báo \\
SQL & Structured Query Language - ngôn ngữ truy vấn cơ sở dữ liệu \\
XGBoost & Extreme Gradient Boosting - thuật toán cây quyết định tăng cường \\
\end{longtable}
